\section{Introducción}
\label{iv:introduccion}
\begingroup
\fontsize{10.5}{12.8}\selectfont
\setlength{\parskip}{1pt plus .5pt minus .25pt}

¿Cómo puede una misma estructura discreta determinar coordenadas del
continuo, sostener una realización excepcional y conservar una dualidad
al pasar a representaciones dimensionales? La cuestión exige algo más
que reunir valores, dimensiones y grupos: exige construir los objetos
intermedios y las aplicaciones que transmiten su información. Este
artículo desarrolla esa conexión en la Holografía Modular Triádica (HMT).
Su resultado articulador es una doble prolongación del estado
genealógico: la incidencia conduce a un álgebra de operadores de vértice
isomorfa a \(V^\natural\), mientras las representaciones de fase, acción
y canales determinan una dualidad compacta y un morfismo de pantallas.
La coordinación de ambas prolongaciones precisa después el dominio
algebraico y las condiciones de transporte hacia una realización de
teoría M.

La construcción comienza antes de la evaluación real.
\emph{All Possible Paths} (APP) designa los caminos sobre el soporte
\(X_9=(\mathbb Z/9\mathbb Z)^2\). Su grafo de Cayley utiliza la ley
aditiva y los generadores cardinales \(\pm(1,0),\pm(0,1)\). Sobre las
mismas posiciones se evalúan suma y producto; la reducción residual
conserva el cociente de cada evaluación. Las hojas aditiva y
multiplicativa son, por tanto, dos evaluaciones de un soporte común.
La firma \(\mathrm{TRIT}=\{-1,0,+1\}\) distingue la orientación del
transporte y los regímenes hiperbólico, parabólico y elíptico de su
realización cuadrática. El \emph{Topological Phase Kernel} (TPK) es la
composición efectiva de selección de evaluación, transporte cardinal
y actualización de registros que prolonga estados compatibles. Estos
estados retienen posiciones, hojas, orientación, residuos, cocientes,
acarreos, prefijos, frontera y memoria. La selección no sustituye al
desplazamiento, y el retorno de una fase no borra la historia recorrida.

% BEGIN R27_FRAMING_INTRO_ES
Las operaciones que conectan el soporte aritmético con la incidencia
admiten pruebas intermedias explícitas. El cambio de módulo, la repetición
del soporte con módulo fijo y la variación dimensional tienen sus propias
leyes del defecto suma--producto. La compresión en tres clases produce
una subretícula de índice dos y una unidad cuadrática de norma uno.
En el emisor, el alfabeto decimal, los empalmes de semiemisiones y el
contraste entre histograma y memoria ordenada conservan las diferencias
entre valor, incidencia y trayectoria. La elección del observable
multiplicativo se examina antes de la selección regional.

% BEGIN R28_FRAMING_INTRO_ES
El multigrafo etiquetado del emisor y su cociente de fase sitúan estas
realizaciones sobre incidencias explícitas. Los cinco cierres mínimos
de ocho aristas cubren los anclajes APP, de propagación, de autoescala
y de cada región de clausura. La ocupación mínima de la frontera y la
paridad de sus lectores de memoria precisan, a su vez, la selección de
representantes y la transformación de los registros al invertir la ruta.
% END R28_FRAMING_INTRO_ES

El primer resultado es la organización conjunta de la estructura
discreta del continuo. En la sección~\ref{iv:continuo-conjunto}, un
mismo árbol de ejecuciones determina la medida cilíndrica y el transporte
de hojas, el balance orientado, la incidencia ternaria, el complejo
de memoria y su ensamblaje terminal con línea determinante. Los mapas
se definen sobre generadores y se prueba su compatibilidad con las
restricciones de historia y las reducciones modulares. La medida
utiliza pesos positivos compatibles en un árbol de ramificación finita
sin hojas terminales; la identidad de Gram conserva el residuo terminal
de supervivencia; la construcción determinante mantiene sus grados y
bases, con las hipótesis adicionales que requiere una evaluación
escalar. Las cinco construcciones conservan así un mismo origen y las
condiciones particulares de cada operación. No son cinco dominios
seleccionados después de producir las coordenadas.

Las secciones de cilindros compatibles determinan las coordenadas
numéricas de la generación correlacionada de \(\pi,\varphi,e,\alpha\)
bajo las condiciones de anidamiento, contracción y frontera
de su realización arquimediana. La proyección excepcional conserva,
en cambio, los datos de incidencia que se realizan como códigos,
retículos y álgebras graduadas. Las hojas areal y volumétrica y sus
medidas, desarrolladas en la sección~\ref{sec:medida-retorno-hojas},
y las vacancias asociadas al refinamiento, en la
sección~\ref{vac-capacidad-trit}, intervienen junto con los registros
regionales en la construcción de \(\alpha\) y en la prolongación
excepcional. La masa electrónica utiliza otra composición de esas
operaciones comunes; su evaluación no es el objeto de este artículo.

La coordenatización que agrupa dos trits en un dígito nonádico dispone de inversa
y de una suma transportada con acarreo; la proyección del cubo ambiente
de \(729\) palabras sobre sus dos primeras coordenadas tiene fibras
uniformes de nueve elementos. La prueba de correlación cuadrática
determina después la identidad de la coordenatización de seis posiciones.
Sobre la frontera, la saturación deja ocho posibilidades, la neutralidad
dual deja dos y la orientación selecciona el representante fijado.
Los controles de clausura conservan los extremos de la identidad
telescópica. 

La conexión excepcional tiene un antecedente concreto en el registro
dodecafásico \(K\). Sus doce componentes proceden del registro firmado
y de las operaciones del TPK. La codificación racional conserva los
bloques cuando se fijan base, longitud, origen de fase y orientación;
la inversión de Hadamard recupera la coordenada firmada correspondiente.
Los soportes del registro intervienen en una bandera marcada del diseño
de Witt. A partir del código ternario se construye el pegado del retículo
de Niemeier de raíces \(A_2^{12}\); la bandera y el origen, junto con la
métrica y la cámara positiva declaradas, seleccionan el vecino sin
raíces que se identifica con el retículo de Leech. Cada operación
transmite un dato necesario al objeto siguiente.

% BEGIN S27_FRAMING
La sección~\ref{star27:comparacion} precisa el paso entre los registros
regionales y la incidencia mediante la comparación de tres estrellas.
Las intersecciones de \(P_\pi\) con \(H_e\), \(H_\varphi\) y
\(H_{\rm APP}\) determinan las cuaternas \(B_K\), \(B_\varphi\) y
\(B_{\rm APP}\); la partición de signo de la cocadena anterior al
acarreo determina la polarización de sus completados. El censo de
las \(495\) cuaternas sitúa estas configuraciones en el diseño completo,
y el estabilizador del marco ordenado de cuatro soportes caracteriza
su rigidez dentro de \(G_W\). Esta rigidez concierne a la acción sobre
posiciones; la determinación numérica de \(K\) conserva su construcción
integral, el origen de fase y la orientación previamente fijados.

La sección~\ref{carry27:seccion} desarrolla el componente entero que
acompaña a la lectura residual. La división euclídea produce los
residuos \(q,a\) y los cocientes \(c,h\), cuya composición conserva la
memoria de los múltiplos de tres. El cociclo, las representaciones
polinómicas sobre \(\mathbb F_3\) y las identidades de transporte
explicitan qué datos atraviesan cada cambio de registro. Su enlace con
los cilindros compatibles conserva las cotas de frontera y las fibras
de la lectura regional. La comparación de estrellas y la ley de
acarreo cumplen así funciones incidenciales y numéricas determinadas
en la conexión que este artículo prolonga hacia retículos y álgebras.

% END S27_FRAMING


Sobre este retículo se construye el álgebra reticular, con osciladores,
extensión central y producto de vértice. La extensión de orden dos de
Frenkel--Lepowsky--Meurman y la extensión de orden tres asociada a la
isometría ternaria realizan \(V^\natural\), una vez comprobadas las
hipótesis de los teoremas de reconocimiento utilizados
\cite{flm,chenlamshimakura2016,abelamyamada2017}.
Los isomorfismos elegidos entre las dos presentaciones conservan vacío,
vector conforme, graduación y productos, y transportan automorfismos
y trazas. El Monstruo es el grupo de automorfismos de esta álgebra;
su acción tiene dominio en \(V^\natural\), no en las cifras de
\(\alpha\). El teorema de Moonshine de Borcherds determina las trazas
graduadas correspondientes~\cite{borcherds}. La construcción del
retículo, la extensión algebraica y el reconocimiento de esas trazas
son etapas distintas de una misma cadena de aplicaciones.

Al alcanzar las trazas de nivel tres, la factorización formal
de la identidad racional determina una divisibilidad válida en todos
los grados, con su exponente óptimo.
% END R27_FRAMING_INTRO_ES

El registro \(K\) posee además una realización direccional.
En la representación declarada de \(A_5\) sobre las doce posiciones,
su proyección sobre el sector tridimensional determina un vector
\(u_K\ne0\) dentro de \(V_{11}=\mathbf1^\perp\).
El vector no nulo determina la recta \(\ell_K\), y su norma cuadrática
exacta normaliza el término de rango uno del proyector sobre
\(V_{10}=V_{11}\cap u_K^\perp\), con núcleo y covariancia
especificados. La reducción \(12\to11\) elimina el modo uniforme;
la reducción \(11\to10\) elimina la dirección seleccionada por el
registro completo. Por otra parte, su órbita cíclica es una referencia
invertible para identificar operadores mediante sus respuestas.
La sección~\ref{sec:acoplamiento-recuperacion} establece las cotas
de estabilidad y el balance bilateral: dos canales construidos a partir de isometrías permiten
recuperar el estado; la extensión unitaria incluye memoria entrante,
y la conservación de un observable requiere la condición de conmutación
de su enunciado.

La realización hilbertiana proporciona el soporte operatorio de la
dualidad. A cada profundidad, las fases etiquetan una base ortonormal
y los operadores de desplazamiento y carácter generan toda su álgebra
matricial. Su relación de Weyl representa la extensión central de las
traslaciones con el signo de la forma de área orientada fijado.
La inclusión mediante la identidad de toda la fibra nueva conserva
unidad y traza normalizada: el límite es el álgebra uniformemente
hiperfinita nonádica y su cierre tracial es un factor de tipo \(II_1\).
Fijar una continuación produce, en cambio, una inclusión en una esquina
y otro cierre. La elección del mapa de refinamiento determina, por
tanto, el tipo de realización y no sólo su dimensión.

La sección de acción y las coordenadas angulares fijan el radio que
especializa la dualidad compacta. La norma determinantal, los
coeficientes del carácter de acción y la renovación regional normalizada
determinan las secciones positivas anterior y posterior al retorno.
Las coordenatizaciones angulares conservan la orientación y el número de vueltas.
Su composición define una elipse etiquetada cuyo producto de semiejes
conserva la acción y cuyo cociente determina un radio adimensional
único. Las condiciones de positividad sitúan las coordenadas generadas
en la coordenatización elíptica; el intercambio de semiejes produce el radio
recíproco conservando la escala común de acción.

Sobre \(\mathbb Z^2\times\mathbb R_{>0}\), el intercambio de las
coordenadas enteras y la inversión del radio forman una involución
que conserva la forma cuadrática compacta. El Hamiltoniano diagonal es
autoadjunto, no negativo y de resolvente compacto; la permutación de
coordenadas lo conjuga unitariamente con el Hamiltoniano de radio
recíproco, incluidos los dominios de operadores y formas cuadráticas.
El cambio de sección residual transporta también la relación \(L_9\)
que conserva el entero. La energía mínima de cada clase está bien
definida y se evalúa mediante dos candidatos enteros. La normalización
por la escala de cuerda expresa la misma dualidad en radios
dimensionales. En el eje imaginario del semiplano superior, la inversión
modular se relaciona con ella mediante una semiconjugación explícita
de parámetros. Este mapa no identifica el operador de intercambio con
los automorfismos de los sectores de un orbifold.

Las pantallas dimensionales coordinan la dirección \(K\), la incidencia
y el sector compacto. La perforación del código ternario produce un
código de longitud once con distancia cinco y partición perfecta.
La reducción ortogonal \(12\to11\to10\) y la reducción reticular
isotrópica \(26\to24\) conservan sus propias presentaciones, relaciones
y núcleos. El morfismo \(\mathcal R_M^{\mathrm{alg}}\) identifica sus
presentaciones cocientadas y entrelaza la dualidad; el grafo de la
proyección conserva la coordenada de \(V_{11}\) junto con su imagen
en \(V_{10}\). Las mismas operaciones se restringen a la colección de
los dos radios determinados por la elipse.

La prolongación dinámica comienza con la composición de trayectorias,
no con la elección de campos. Al conservar el dato precedente en la
frontera, los cambios de dirección y de tipo aritmético definen un
cociclo aditivo de acción. La superficie de mundo y la acción de
Polyakov especifican después una realización variacional. El mapa
hacia la realización de teoría M recibe las pantallas en un codominio
undecadimensional con métrica, tres-forma, gravitino, graduación,
acción y restricciones. Bajo las hipótesis de entrelazamiento y
transporte de dominios, el teorema de supercarga conserva \(Q^2=H\)
sobre la imagen de la isometría. La identificación física completa
requiere además las tensiones, los osciladores y amplitudes, el control
de anomalías y los límites declarados en la
sección~\ref{iv:condiciones-realizacion-fisica}.

Esta jerarquía determina el desarrollo: núcleo y continuo conjunto;
generación coinductiva, registro \(K\) y las dos construcciones de
\(\alpha\); incidencia y álgebra excepcional; realización hilbertiana,
acción y radio; dualidad, pantallas y transporte dinámico.
La sección~\ref{sec:k-registro} determina el registro utilizado en
\ref{exc:k-direccion}; la realización de \ref{iv:hilbert} sostiene
los operadores de \ref{iv:dualidad}; y
\ref{iv:pantallas} compone las reducciones en el morfismo que recibe la
realización undecadimensional. El argumento enlaza así estructura,
valor, representación y ecuación mediante sus aplicaciones efectivas.

\par\endgroup
